\documentclass{../tex/ficheseance}
% ON TESTE ICI SANS LE \usepackage{babel} POUR VERIFIER LE FALLBACK FRANCAIS

% --- INCLUSION DU CATALOGUE (Optionnel, peut être dans un fichier à part) ---
\DefinirSituation{theme01-obj03}{
	objectif = {Apprendre à viser avec un viseur},
	consignes = {
		\begin{itemize}
			\item Regarder la cible au travers du viseur
			\item Garder les deux yeux ouverts
			\item Ne pas crisper la main d'arc
		\end{itemize}
	},
	criteres = {Groupement dans le jaune à 10m}
}

\DefinirSituation{echauffement}{
	objectif = {Échauffement articulaire et musculaire},
	temps = {15 min},
	situations = {En cercle sans matériel},
	consignes = {Réveil articulaire complet (nuque, épaules, poignets)}
}
% ----------------------------------------------------------------------------

\begin{document}
% Informations de la séance
\numfiche{3}
\datefiche{23 Octobre 2026}
\theme{Visée et Continuité}
\objectifseance{Intégrer l'usage du viseur pour les débutants.}
\archer{Groupe Initiation}

% Génération de l'en-tête
\makeficheheader

% Création du tableau de séance
\begin{tableauSeance}

	% Exemple 1 : Utilisation directe d'une situation du catalogue sans modification
	\ligneDepuisCatalogue{echauffement}
	
	% Exemple 2 : Surcharge de valeurs (on ajoute ce qui manque pour cette séance précise)
	\ligneDepuisCatalogue[
		temps = {30 min},
		situations = {Tir sur blason de 80cm à 10m},
		outils = {Arc avec viseur réglé},
		evaluations = {Comptage des flèches dans le jaune}
	]{theme01-obj03}
	
	% Exemple 3 : Surcharge pour modifier une valeur existante du catalogue (on change les consignes pour CE groupe)
	\ligneDepuisCatalogue[
		temps = {15 min},
		situations = {Tir de confirmation},
		consignes = {\textbf{Rappel important :} bien insister sur le placement des doigts !} % Ceci écrase le catalogue
	]{theme01-obj03}

	% Exemple 4 : Utilisation de la commande manuelle standard (hors catalogue)
	\ligneSeance{
		Retour au calme
	}{
		10 min
	}{
		Sur le pas de tir
	}{
		Aucun
	}{
		Rangement du matériel, étirements légers
	}{
		Le matériel est rangé correctement
	}{
		Feedback oral
	}

\end{tableauSeance}

\end{document}
